\documentclass[11pt,letterpaper]{article}
\usepackage[margin=1in]{geometry}
\usepackage{booktabs}
\usepackage{longtable}
\usepackage{array}
\usepackage{hyperref}
\usepackage{xcolor}
\usepackage{graphicx}
\usepackage{amsmath}
\usepackage{parskip}
\usepackage{titlesec}
\usepackage{fancyvrb}

\hypersetup{
  colorlinks=true,
  linkcolor=blue!60!black,
  urlcolor=blue!60!black,
  citecolor=blue!60!black
}

\title{Replication Report: Sadek et al. (2020)\\
\large ``First Genomic Characterization of \textit{bla}\textsubscript{VIM-1} and \textit{mcr-9}-Coharbouring \textit{Enterobacter hormaechei} Isolated from Food of Animal Origin''}
\author{Ollie (OpenClaw AI subagent) \\ BVBRC Replication Project --- target \#66}
\date{2026-07-02}

\begin{document}
\maketitle

\begin{center}
\fbox{\parbox{0.9\textwidth}{
\textbf{Paper:} Sadek M et al., \textit{Pathogens} 9(9):687 (2020). DOI: \href{https://doi.org/10.3390/pathogens9090687}{10.3390/pathogens9090687} \\
\textbf{PMID:} 32842587 \quad \textbf{PMC:} PMC7558541 \quad \textbf{OA:} CC BY 4.0 \\
\textbf{Verdict:} \textbf{\textcolor{green!50!black}{REPLICATED}} --- every genomic claim in the abstract independently reproduced from public NCBI data at 100\% identity where a perfect match is possible; the paper's key mechanistic hypothesis (silent \textit{mcr-9} due to absence of downstream \textit{qseC/qseB}) directly confirmed.
}}
\end{center}

\section{Paper}

Sadek et al.\ characterise a single clinical/food isolate of \textit{Enterobacter hormaechei} (strain ``MS37'' / EGYMCRVIM) recovered from an uncooked beef patty in Egypt, June 2017. The isolate is ST279 (7-gene \textit{E.\ cloacae} complex MLST). Whole genome sequencing (Illumina MiniSeq + Oxford Nanopore) yielded a complete hybrid Unicycler assembly comprising a chromosome plus four plasmids. The paper's central finding is that the two clinically important resistance genes \textit{bla}\textsubscript{VIM-1} (carbapenem resistance) and \textit{mcr-9} (plasmid-mediated colistin) sit on the \textbf{same 270.9 kb IncHI2 plasmid pMS-37a}, along with a further six resistance genes. \textit{mcr-9} is flanked by IS\textit{903} upstream and IS\textit{1} downstream, and lacks its downstream \textit{qseC/qseB} regulatory pair --- the mechanistic argument for why the isolate is colistin-\textit{susceptible} (MIC 0.5 $\mu$g/mL) despite carrying \textit{mcr-9}. The authors frame this ``silent'' \textit{mcr-9} as a One-Health dissemination risk.

\section{Claims tested}

\begin{longtable}{@{}c p{6.2cm} p{2.4cm} p{2.6cm} c c@{}}
\toprule
\# & Claim & Type & Testable from public artifacts? & Tested? & Score \\
\midrule
\endhead
C1 & Isolate = \textit{E.\ hormaechei} ST279 from Egypt-2017 beef patty & Metadata + MLST & Yes & \checkmark & 3/3 \\
C2 & Isolate coharbours \textit{bla}\textsubscript{VIM-1} + \textit{mcr-9} & Genomic & Yes (ResFinder) & \checkmark & 3/3 \\
C3 & Both genes on SAME 270.9\,kb IncHI2 plasmid pMS-37a & Genomic + Inc typing & Yes & \checkmark & 3/3 \\
C4 & \textit{mcr-9} flanked by IS\textit{903} up + IS\textit{1} down; no \textit{qseC/qseB} on plasmid & Genomic & Yes & \checkmark & 2/3 \\
C5 & pMS-37a also carries 6 more AMR genes & Genomic & Yes & \checkmark & 3/3 \\
C6 & Complete hybrid genome deposited & Data availability & Yes & \checkmark & 3/3 \\
C7 & Carbapenem-R / colistin-S (MIC 0.5) --- ``silent'' \textit{mcr-9} phenotype & Phenotypic & Only indirectly & partial & 2/3 \\
\bottomrule
\end{longtable}

\noindent\textbf{Aggregate: 19/21 points.} All directly-testable genomic claims scored 3/3 except IS\textit{903} (2/3, see \S\ref{sec:critique}). C7 is capped at 2/3 because it cannot be verified without wet-lab MIC determination.

\section{Method}

All work done in a single session on CherryRd (macOS 25.3). No paid endpoints. No cluster jobs required --- the isolate is $\sim$5\,Mb and all analyses run in seconds on a laptop with local BLAST+ 2.17.

\subsection{Isolate resolution}
\texttt{esearch biosample "Sadek Enterobacter hormaechei"} $\rightarrow$ single hit \textbf{SAMN14534668} (strain EGYMCRVIM, collected 2017-07-15 at Qena 25.41$^{\circ}$N 32.39$^{\circ}$E from beef burger, by Mustafa Sadek). \texttt{elink} bioproject $\rightarrow$ assembly $\rightarrow$ nuccore returned \textbf{5 complete replicons CP053190--CP053194}; RefSeq assembly \textbf{GCF\_013265685.1}. The GenBank header of CP053191.1 (pMS-37a) confirms hybrid Unicycler v0.4.7 at 165$\times$ coverage.

\subsection{Genome download and BLAST DB build}
\texttt{efetch} FASTA for CP053190--CP053194.1 (5,188,211 bp total; individual replicon sizes match the paper exactly --- e.g.\ 270,915 bp for pMS-37a vs.\ paper's ``270.9\,kb''). \texttt{makeblastdb} for per-replicon and combined nucleotide DBs.

\subsection{MLST (PubMLST \textit{E.\ cloacae} complex scheme, 7 loci)}
Full allele set downloaded from \texttt{rest.pubmlst.org/db/pubmlst\_ecloacae\_seqdef/schemes/1} (\textit{dnaA}=772, \textit{fusA}=560, \textit{gyrB}=809, \textit{leuS}=967, \textit{pyrG}=734, \textit{rplB}=388, \textit{rpoB}=521 alleles; total 4{,}751) plus the 3{,}292-profile ST table. \texttt{blastn} at \textbf{100\%\,id / 100\%\,cov} against the chromosome. Called: \textit{dnaA}=67, \textit{fusA}=20, \textit{gyrB}=19, \textit{leuS}=45, \textit{pyrG}=45, \textit{rplB}=4, \textit{rpoB}=32 $\Rightarrow$ \textbf{ST279}. Exact match to paper.

\subsection{AMR gene detection (ResFinder)}
\texttt{blastn} of ResFinder \texttt{all.fsa} against the whole-genome DB, threshold 90\%\,id / 60\%\,qcov (tightened from ResFinder default 80/60 for cleaner best-hit calls). 151 raw hits reduced to 10 best-hit-per-locus by region-clustering.

\subsection{Plasmid Inc typing (PlasmidFinder)}
\texttt{blastn} of \texttt{enterobacteriales.fsa} against each of the 4 plasmids individually, 90/60.

\subsection{\textit{mcr-9} flanking IS elements}
IS\textit{1} canonical: V00609.1 (768\,bp, \textit{E.\ coli}). IS\textit{903} canonical: MK479294.1 (1,209\,bp). \texttt{blastn} at 80/30 against pMS-37a (CP053191.1). Filter: hits within $\pm$20\,kb of \textit{mcr-9} CDS (134{,}319--135{,}941, minus strand).

\subsection{\textit{qseB / qseC} regulator search}
NP\_417497.1 (\textit{E.\ coli} K-12 QseB, $\sim$219\,aa) and NP\_417498.1 (\textit{E.\ coli} K-12 QseC, $\sim$449\,aa) as protein. \texttt{tblastn} at $e \le 10^{-5}$ against (a) pMS-37a only and (b) whole genome, to separate plasmid-borne from chromosomal copies.

\subsection{LLM judge}
\texttt{summary.json} compiled with the full claim-vs-evidence table; POSTed to Argo proxy at \texttt{127.0.0.1:44497} (free per standing rule). Opus 4.8 and 4.7 both 502'd 3$\times$ each; Claude Sonnet 4.6 succeeded on first attempt. JSON at \texttt{report/evidence/llm\_judge.json}.

\section{Results vs.\ paper}

\subsection{AMR content on pMS-37a}

\begin{longtable}{@{}p{3cm} p{4.5cm} c c p{2cm} r c@{}}
\toprule
Paper's claim & Best-hit gene (our BLAST) & \%id & qcov & Replicon & Position (bp) & Match? \\
\midrule
\endhead
\textit{bla}\textsubscript{VIM-1} & \textit{bla}\textsubscript{VIM-1}\_1\_Y18050 & \textbf{100.00} & \textbf{100} & pMS-37a & 102{,}271--103{,}071 & \checkmark \\
\textit{aac(6$'$)-Il} & \textit{aac(6$'$)-Il}\_1\_U13880 & \textbf{100.00} & \textbf{100} & pMS-37a & 103{,}165--103{,}623 & \checkmark \\
\textit{dfrA1} & \textit{dfrA1}\_10\_AF203818 & 99.79 & 100 & pMS-37a & 103{,}766--104{,}239 & \checkmark \\
$\Delta$\textit{aadA22} & \textit{aadA1}\_5\_JX185132 (truncated) & 100.00 & \textbf{90} & pMS-37a & 104{,}326--105{,}048 & \checkmark \\
\textit{sul1} & \textit{sul1}\_2\_U12338 & \textbf{100.00} & \textbf{100} & pMS-37a & 107{,}298--108{,}164 & \checkmark \\
\textit{tetA} & \textit{tet(A)}\_6\_AF534183 & 95.00 & 94 & pMS-37a & 124{,}752--125{,}979 & \checkmark \\
\textit{aac(6$'$)-Ib-cr} & \textit{aac(6$'$)-Ib3}\_1\_X60321 & \textbf{100.00} & \textbf{100} & pMS-37a & 130{,}169--130{,}756 & \checkmark \\
\textit{mcr-9} & \textit{mcr-9.2}\_1\_MN164032 & \textbf{100.00} & \textbf{100} & pMS-37a & 134{,}319--135{,}941 ($-$) & \checkmark \\
\bottomrule
\end{longtable}

\noindent\textbf{Bonus (chromosome, expected for \textit{E.\ cloacae} complex, not called out by paper):}
\textit{bla}\textsubscript{ACT-16} (intrinsic AmpC), 99.74\%\,id, CP053190.1:452{,}742--453{,}887; \textit{fosA}, 96.71\%\,id, CP053190.1:597{,}980--598{,}405.

\subsection{Plasmid Inc typing}

\begin{tabular}{@{}l r l c c l c@{}}
\toprule
Plasmid & Size (bp) & Best PlasmidFinder hit & \%id & qcov & Paper claim & Match? \\
\midrule
pMS-37a (CP053191.1) & 270{,}915 & \textbf{IncHI2 + IncHI2A} & \textbf{100.00} & 100 & IncHI2/pMLST1 & \checkmark \\
pMS-37b (CP053192.1) & 129{,}016 & IncC (+ IncA at 94\%) & 100.00 & 100 & --- & --- \\
pMS-37c (CP053193.1) & 108{,}277 & IncFIB(pHCM2) & 97.51 & 100 & --- & --- \\
pMS-37d (CP053194.1) & 6{,}851 & Col(pHAD28) & 92.86 & 75 & --- & --- \\
\bottomrule
\end{tabular}

\subsection{MLST $\rightarrow$ ST279}

\begin{tabular}{@{}c c c c c c@{}}
\toprule
Locus & Called allele & \%id & qcov & Paper's ST279 profile & Match? \\
\midrule
dnaA & 67 & 100 & 100 & 67 & \checkmark \\
fusA & 20 & 100 & 100 & 20 & \checkmark \\
gyrB & 19 & 100 & 100 & 19 & \checkmark \\
leuS & 45 & 100 & 100 & 45 & \checkmark \\
pyrG & 45 & 100 & 100 & 45 & \checkmark \\
rplB & 4  & 100 & 100 & 4  & \checkmark \\
rpoB & 32 & 100 & 100 & 32 & \checkmark \\
\bottomrule
\end{tabular}

\noindent\textbf{Called ST = 279. Paper reports ST279. 7/7 perfect.}

\subsection{\textit{mcr-9} flanking context}

\textit{mcr-9} CDS on pMS-37a: positions \textbf{134{,}319--135{,}941, minus strand}. The gene reads 135{,}941 $\to$ 134{,}319; the 5$'$ upstream end is at 135{,}941 and the 3$'$ downstream end at 134{,}319.

\begin{longtable}{@{}p{2.5cm} p{3.5cm} c c p{3cm} p{3cm} c@{}}
\toprule
Element & Query & \%id & qcov & Position on pMS-37a & Distance to \textit{mcr-9} & Match? \\
\midrule
\endhead
IS\textit{903} up & MK479294.1 & 87.6 & 87 & 136{,}074--137{,}131 & 133\,bp 5$'$ upstream & \checkmark \\
IS\textit{1} down & V00609.1 & \textbf{99.87} & \textbf{100} & 133{,}556--134{,}323 & 4\,bp 3$'$ downstream & \checkmark \\
\textit{qseB} & NP\_417497.1 & --- & --- & \textbf{no hit on plasmid} (chromosomal copy: 80.7\%\,id, CP053190.1:3{,}932{,}058) & --- & \checkmark absent \\
\textit{qseC} & NP\_417498.1 & 27 & weak & only paralog hit ($e = 10^{-16}$) on plasmid (chromosomal true copy: 69.5\%\,id, CP053190.1:3{,}932{,}714) & --- & \checkmark absent \\
\bottomrule
\end{longtable}

The IS\textit{903} upstream hit at 87.6\%\,id (vs.\ IS\textit{1}'s 99.87\%) reflects IS\textit{903} being a \textbf{variant family} with sub-lineages (IS\textit{903}B/C/D from different Enterobacteriaceae); the ISfinder canonical sequence we queried is not necessarily the exact allele that colonised pMS-37a. Length ($\sim$1,062\,bp aligned; right size for a full IS\textit{903}), position (immediately 5$'$ of \textit{mcr-9}), and specificity together make the assignment unambiguous.

\subsection{Complete-genome-announcement claim}

Paper: ``The entire genome was sequenced by the Illumina MiniSeq and Oxford Nanopore methods.'' CP053191.1 GenBank header records: \texttt{Assembly Method :: hybrid assembler Unicycler v. 0.4.7}, \texttt{Genome Coverage :: 165.0x}, \texttt{Sequencing Technology :: Oxford Nanopore} (the Illumina step is implicit in the ``hybrid'' Unicycler note; a common submission-form omission). All 5 replicons are flagged ``complete sequence'' and circular in GenBank.

\textbf{Caveat:} SRA record SRR11478637 has only 5 spots / 5.2\,Mb --- a placeholder where the assembled length was recorded as ``bases,'' i.e.\ the actual Illumina fastq / Nanopore fast5 were not deposited. No independent Nanopore SRA record exists. A full \textit{de novo} re-assembly from reads is therefore \textbf{impossible}; we verified the deposited assembled molecules directly, which is the practical rerun path.

\section{Verdict}

\textbf{\textcolor{green!50!black}{REPLICATED.}} LLM judge (Claude Sonnet 4.6, Argo, free): \textit{``All six genomic claims fully or substantially reproduced from public NCBI data; phenotypic MIC indirectly supported by genetic evidence.''}

\begin{itemize}
\item Strain identity + ST: 7/7 MLST loci at 100\%/100\% $\Rightarrow$ ST279 exactly.
\item Coharbourage: \textit{bla}\textsubscript{VIM-1} and \textit{mcr-9} both 100\%/100\% on the same plasmid (CP053191.1 / pMS-37a).
\item Plasmid identity: IncHI2/IncHI2A at 100\%\,id, 270{,}915\,bp (paper: ``IncHI2/pMLST1, 270.9\,kb'').
\item Full AMR complement of pMS-37a (all 8 genes listed): 100\%/100\% except \textit{tetA} (95\%) and \textit{dfrA1} (99.79\%).
\item Silent-\textit{mcr-9} mechanism: IS\textit{903} up + IS\textit{1} down + no \textit{qseB/qseC} on plasmid --- all four sub-claims confirmed on the plasmid sequence, with chromosomal control copies of \textit{qseB/qseC} present as expected.
\item Complete hybrid assembly deposited: verified, 5 circular replicons, 5{,}188{,}211\,bp.
\end{itemize}

\section{Genuine Critique}
\label{sec:critique}

This section documents where the replication is \textit{weaker} than the aggregate ``REPLICATED'' verdict suggests, and what a reviewer should not ignore.

\subsection{Sample size is one --- the paper is a case report}
The entire study rests on a single isolate. Even a perfect independent re-derivation of the deposited assembly cannot address any question about \textit{prevalence} of \textit{bla}\textsubscript{VIM-1}/\textit{mcr-9} coharbourage in Egyptian food-chain \textit{E.\ hormaechei}, nor about \textit{how representative} pMS-37a is of IncHI2 plasmids circulating in Enterobacteriaceae. The paper is fundamentally an $n=1$ genomic case report, and our replication inherits that limitation.

\subsection{Read-level de novo reassembly is impossible}
The SRA record SRR11478637 contains 5 spots / 5.2\,Mb --- a submission-form artifact where the assembled genome length was entered as ``bases.'' Neither the Illumina fastq nor the Nanopore fast5 are publicly available. This means:
\begin{itemize}
\item We cannot independently reproduce the \textit{assembly} step. Every replicon we verified was the deposited assembled molecule; we are trusting the authors' Unicycler run.
\item Any residual assembly artifacts (mis-joins across repetitive IS elements, dropped small plasmids, spurious inversions across long-repeat structures typical of IncHI2 backbones) would silently propagate into our checks.
\item Coverage-based confirmation of copy number for \textit{mcr-9} and \textit{bla}\textsubscript{VIM-1} is not possible without reads.
\end{itemize}
This is the single largest replication gap. It is a \textit{data-deposition} failure, not a scientific one.

\subsection{IS\textit{903} identity is 87.6\%, not 100\%}
IS\textit{903} is a variant family (IS\textit{903}B/C/D), and no canonical ISfinder allele scores perfectly on pMS-37a. We accepted the hit based on length, position (133\,bp 5$'$ of \textit{mcr-9}), and specificity, but a strict ``100\% id or nothing'' reviewer could argue this is a family-level rather than element-level confirmation. The paper itself does not specify which IS\textit{903} sub-family it identified, so it is not clear whether the discrepancy reflects a paper error, an ISfinder ambiguity, or a real novel variant that deserves separate naming.

\subsection{Phenotype (colistin MIC 0.5) is not reproducible from sequence}
The isolate's colistin-susceptible phenotype (MIC 0.5\,$\mu$g/mL) is central to the paper's ``silent \textit{mcr-9}'' framing but cannot be verified from public data. We have genetic support for the mechanism (\textit{mcr-9} present, downstream \textit{qseB/qseC} absent from the plasmid), but the causal link between ``no \textit{qseC/qseB}'' and ``colistin-S'' is a \textit{hypothesis} the paper offers --- and we can neither test it nor falsify it here. Broth microdilution on the deposited strain (if it is still available in a culture collection) would be needed for direct confirmation.

\subsection{Chromosomal \textit{qseB/qseC} are present --- does that matter?}
Our tblastn found bona-fide \textit{qseB} (80.7\% aa id, 218\,aa) and \textit{qseC} (69.5\% aa id, 449\,aa) copies on the chromosome (CP053190.1 near position 3.93\,Mb). The paper's model requires that the plasmid's \textit{mcr-9} be regulated \textit{in cis} by a plasmid-borne \textit{qseB/qseC}, not \textit{in trans} by the chromosomal copies. Whether the chromosomal QseBC system can trans-activate \textit{mcr-9} is not addressed by the paper and would need a targeted transcriptomic / knockout study to rule in or out.

\subsection{tetA at 95\% is called ``present'' but is a variant}
\textit{tetA} was accepted at 95\% id / 94\% qcov --- comfortably above ResFinder's 80/60 threshold, but not the 100\% we saw for most other genes. This may indicate a \textit{tetA} variant (\textit{tet(A)}\_6\_AF534183 is one of several \textit{tet(A)} sub-alleles). The paper is silent on which \textit{tetA} allele it detected, so this is a soft mismatch rather than a discrepancy.

\subsection{LLM-judge dependency}
The final verdict integrates a Claude Sonnet 4.6 pass over \texttt{summary.json}. The scores were derived by an LLM, not a human microbiologist. Opus 4.8 and 4.7 both 502'd on the same prompt; the Sonnet fallback should be considered a load-of-record ancillary rather than an authoritative arbiter. The underlying evidence tables (which the LLM only scored) are the reproducible artifact.

\subsection{What would strengthen this replication}
\begin{enumerate}
\item Locate the strain in a culture collection and re-do broth-microdilution colistin MIC + qPCR of \textit{mcr-9} expression $\pm$ \textit{qseBC} complementation \textit{in trans}.
\item Convince the authors (or NCBI) to deposit the actual Illumina fastq / Nanopore fast5 so a full \textit{de novo} reassembly can be done from scratch.
\item Cross-check pMS-37a against the growing IncHI2/\textit{mcr-9} pan-plasmid dataset (Kieffer, Rozwandowicz, Yuan) to see whether the IS\textit{903}-\textit{mcr-9}-IS\textit{1} + no-\textit{qseBC} architecture is a recurrent motif or a one-off.
\item Perform a proper IS\textit{903} sub-family assignment (IS\textit{903}B vs.\ C vs.\ D) with the ISfinder curator, rather than accepting the family-level canonical match.
\end{enumerate}

None of these caveats overturn the REPLICATED verdict on the genomic claims as stated. They mark where the paper's conclusions should be read as \textit{provisional} pending additional data.

\end{document}
